\documentclass[a4paper,12pt]{article}
\usepackage[utf8]{inputenc}
\usepackage[spanish]{babel}
\usepackage{color}
\usepackage{parskip}
\usepackage{graphicx}
\usepackage{multirow}
\usepackage{listings}
\usepackage{vmargin}
\graphicspath{ {imagenes/} }
\definecolor{mygreen}{rgb}{0,0.6,0}
\definecolor{lbcolor}{rgb}{0.9,0.9,0.9}
\usepackage{epstopdf}


\setpapersize{A4}
\setmargins{2.5cm}       % margen izquierdo
{1.5cm}                        % margen superior
{16.5cm}                      % anchura del texto
{23.42cm}                    % altura del texto
{10pt}                           % altura de los encabezados
{1cm}                           % espacio entre el texto y los encabezados
{0pt}                             % altura del pie de página
{2cm}     

\lstset{
backgroundcolor=\color{lbcolor},
    tabsize=4,    
%   rulecolor=,
    language=C++,
        basicstyle=\tiny,
        aboveskip={1.5\baselineskip},
        columns=fixed,
        showstringspaces=false,
        extendedchars=false,
        breaklines=true,
        prebreak = \raisebox{0ex}[0ex][0ex]{\ensuremath{\hookleftarrow}},
        frame=single,
        showtabs=false,
        showspaces=false,
        showstringspaces=false,
        identifierstyle=\ttfamily,
        keywordstyle=\color[rgb]{0,0,1},
        commentstyle=\color[rgb]{0.026,0.112,0.095},
        stringstyle=\color{red},
        numberstyle=\color[rgb]{0.205, 0.142, 0.73},
%        \lstdefinestyle{C++}{language=C++,style=numbers}’.
}

\begin{document}

\begin{Large}
 NOMBRE-$>$CHRISTOFER FABIÁN CHÁVEZ CARAZAS \par
 TEMA-$>$ A
\end{Large}

\section{Código}

  \begin{lstlisting}
MODULE main
DEFINE
	n1 := 1;
	n2 := 2;
VAR
	x: {1,2};
	P1: process persona(P2.b_my,x,n1);
	P2: process persona(P1.b_my,x,n2);
LTLSPEC !(F(P1.state = l3 & P2.state = l3));



MODULE persona(b_other,x,number)
VAR
	b_my: boolean;
	state : {l1,l2,l3};
ASSIGN
	init(b_my) := TRUE;
	init(state) := l1;
	next(state) := case
						state = l1 : l2;
						state = l2 & (x = number | b_other = FALSE) : l3;
						state = l3 : l1;
						TRUE: state;
						esac;
	next(b_my) := case
						state = l1 : TRUE;
						state = l3 : FALSE;
						TRUE: b_my;
						esac;
	next(x) := case
					state = l1 & number = 1: 2;
					state = l1 & number = 2: 1;
					TRUE: x;
					esac;

  \end{lstlisting}

\section{Sentencia}

Lo que nos pide el problema es probar que dos personas no puedan ingresar a la cuenta bancaria al mismo tiempo.
Esto se podría traducir en la siguiente sentencia LTL:
$$!(F(P1.state = l3 \ \& \ P2.state = l3))$$
Si un objeto persona esta en el estado l3, entonces esta dentro de la cuenta bancaria: ``En un futuro dos personas no podran estar en la cuenta bancaria al mismo tiempo''

\section{Test}

\begin{figure}[h]
 \centering
 \includegraphics[scale = 0.5]{1.png}
\end{figure}

  

\end{document}
